\section{Descrizione del progetto}
\label{sec:descrizione}

TicketSeller è un programma per vendere biglietti di eventi come concerti, spettacoli teatrali e altri eventi culturali. Lo usano due tipi di persone: i \emph{clienti}, che comprano i biglietti da soli su una postazione self-service chiamata \emph{kiosk}, e i \emph{cassieri}, cioè i dipendenti della biglietteria che gestiscono l'elenco degli eventi.

Per ogni evento si registrano il nome, la data e l'ora, la sala, una descrizione facoltativa, il numero di posti e il prezzo del biglietto. Ogni posto dell'evento ha un suo biglietto, che può essere \emph{disponibile}, \emph{bloccato} (riservato per pochi minuti a un cliente che sta pagando) oppure \emph{venduto}.

Il cliente vede l'elenco degli eventi in programma, ne sceglie uno e vede la mappa dei posti con lo stato e il prezzo di ciascuno. Seleziona i posti liberi che vuole e il programma gli mostra il totale. Poi inserisce nome, e-mail e dati della carta di credito e conferma l'acquisto.

A quel punto il programma controlla che i posti siano ancora liberi e li blocca a nome del cliente, così nessun altro può prenderli. Calcola il totale e chiede il pagamento al servizio della carta di credito (nel progetto il pagamento è simulato). Se il pagamento va a buon fine i biglietti diventano venduti e il cliente riceve un codice che vale come biglietto; se il pagamento non va a buon fine, oppure i dati della carta sono sbagliati, i posti tornano liberi. Un posto non può restare bloccato più di 15 minuti: passato questo tempo torna libero da solo.

Il prezzo può essere calcolato in tre modi: prezzo normale, sconto per i gruppi (sconto in percentuale da un certo numero di biglietti in su) e abbonamento. Un abbonamento ha un codice, vale per alcuni eventi e dà uno sconto: se il cliente inserisce il codice, i biglietti degli eventi compresi nell'abbonamento costano meno, gli altri hanno il prezzo normale.

Il cassiere può aggiungere nuovi eventi, purché in una data futura, modificarli (senza scendere sotto il numero di biglietti già venduti) ed eliminarli, purché nessun biglietto sia stato venduto o sia in fase di pagamento.

I dati vengono salvati sul computer, così restano disponibili anche dopo aver chiuso il programma.
